%\chapter{\rm{VASP}的编译}%完全指南：编译器、数学库、MPI库与GPU编译详解
\chapter{\rm{VASP}的编译与输出}
因为\textrm{VASP}是在\textrm{Linux}系统下运行的，
%（Vienna Ab initio Simulation Package）作为第一性原理计算领域最核心的软件之一，其
用户有必要了解其编译安装的技术细节。%必须面对的技术挑战。
%编译安装。与许多一键安装的普通软件不同，
科学计算软件常有大量源代码文件，用户不大可能逐个对源代码文件进行编译和链接，生成最终的可执行文件。在\textrm{Linux}环境下，大型的科学计算软件的编译和维护，主要依赖\textrm{make}命令工具\footnote{此处默认用户熟悉并理解基本的\textrm{Linux}的操作命令知识。}。\textrm{make}命令通过执行安装文件\textrm{makefile}指定的全部指令、实现软件的编译、链接的自动化构建，提高软件开发的效率。一般地讲，科学计算的过程常会涉及大规模的矩阵运算、快速傅里叶变换\textrm{(Fast Fourier Transformation, FFT)}和并行通信，计算任务的效率也高度依赖编译器优化、数学库的算法实现以及并行通信库的性能。因此，科学计算软件的编译，根源于其对底层计算资源的深度理解和掌握。\textrm{VASP}的编译也不例外，主要就是明确安装文件\textrm{makefile}中编译器版本、数学库、\textrm{MPI (Message Passing Interface)}并行库以及可选的\textrm{GPU}加速支持等变量的设定，以及涉及的编译、链接命令的参数的意义。%任何一个环节的疏漏都可能导致编译失败或运行时性能低下。\textrm{VASP}的编译也不例外。
理解\textrm{VASP}编译的每一个环节，不仅是成功安装\textrm{VASP}的前提，也是后续计算能够高效运行的基础。

%本文将从零开始，
这里介绍\textrm{VASP}编译的完整流程，涵盖编译器的选择、重要数学库(包括\textrm{BLAS}、\textrm{LAPACK}、\textrm{ScaLAPACK}、\textrm{FFTW})的配置、\textrm{MPI}库的配置，设置\textrm{GPU}加速(\textrm{OpenACC}与\textrm{CUDA})的编译选项。重点是安装文件\textrm{makefile.include}的必要说明。%以及常见问题的诊断与解决。无论你是初次接触VASP的新手，还是希望优化现有编译配置的老手，本文都将提供详实的参考。

\section{\rm{VASP}编译的系统要求}
%在开始编译之前，首先需要确认系统是否满足VASP的基本要求。根据VASP官方文档，
%编译\textrm{VASP}需要的组件有
%### 1.1 编译器要求
\textrm{VASP}源代码主要用\textrm{Fortran}语言编写(需要使用\textrm{F2008}及以上标准的编译器)，同时也包含少量的\textrm{C}和\textrm{C++}代码。官方推荐的编译器套件包括：
\begin{itemize}
	\item \textrm{Intel oneAPI}~(含\textrm{ifort}编译器)
	\item \textrm{NVIDIA HPC-SDK}~(含\textrm{nvfortran}编译器)\footnote{特别地，对于\textrm{GPU}版本的\textrm{OpenACC}编译，建议使用\textrm{NVIDIA HPC-SDK}(版本$\geqslant$21.2)}
	\item \textrm{GCC/GFortran}~(版本要求$\geqslant\mathrm{7.X}$，从\textrm{VASP~6.2.0}开始支持)
        \item \textrm{AMD AOCC}~(\textrm{AMD}优化编译器，适用于\textrm{AMD EPYC}处理器)
\end{itemize}
用户可以根据自己的硬件配置，选择符合要求的一款编译器安装即可，类似\textrm{Intel openAPI}这种商业套装，包含了编译器和各类库函数等完整工具，且针对\textrm{X86}架构的硬件有专门优化。而开源的\textrm{GCC/GFortran}等开源编译器，相应的数学库和\textrm{MPI}库都需要独立安装，对用户在操作系统方面的基础知识有一定要求。在高性能计算集群中，相关的编译器一般都是安装、配备好的，用户只需向集群运维和管理员确认编译器环境加载方案，遵照运行即可。用户可以通过类似\texttt{which ifort}等命令，检查相关的软件和环境变量是否已经正确安装或加载。

%### 1.2 必需数值库
\textrm{VASP}编译主要依赖的重要数学函数库包括：
%见下表。关键数学库(\textrm{Math Kernel Library, MKL})是\textrm{intel}的}：
%\begin{table}[h!]
%    \centering
%    \begin{tabular*}{\textwidth}{|@{\extracolsep{\fill}}c|@{\extracolsep{\fill}}c|@{\extracolsep{\fill}}c|}
%        \toprule
% 库名称  &功能  &备注\\ 
%\midrule
%\textrm{BLAS} &基础线性代数运算 &矩阵乘法、向量运算等\\
%\textrm{LAPACK} &高级线性代数运算 &特征值求解、矩阵分解等\\
%\textrm{ScaLAPACK} &并行线性代数运算 &分布式内存并行版本\\
%\textrm{FFTW} &快速\textrm{Fourier}变换 &处理平面波基组的核心库\\
%        \bottomrule
%    \end{tabular*}
%    \caption{\textrm{VASP}编译最依赖的数学库。备注如果使用\textrm{Intel oneAPI}套件，\textrm{BLAS}、\textrm{LAPACK}和\textrm{ScaLAPACK}已包含在重要数学库\textrm{(Math Kernel Library, MKL)}中，无需单独安装。}
%\label{Tab:Compile-Tag-0}
%\end{table}
\begin{itemize}
	\item \textrm{BLAS}:~完成基础线性代数运算，如矩阵乘法、向量运算等。
	\item \textrm{LAPACK}:~在\textrm{BLAS}基础上，定义并完成了更多高级线性代数运算，如矩阵的分解或矩阵特征值求解等。
	\item \textrm{ScaLAPACK}:~完成分布式内存、并行环境下各类线性代数运算。
	\item \textrm{FFTW}:~实现快速\textrm{Fourier}变换的有效工具。
\end{itemize}
\textrm{VASP}计算中会有大量矩阵的构造、分解和矩阵对角化，物理量也会有实空间和倒空间的多次变换，因此这些库函数是完成计算所必须的。一般地，如果用户使用\textrm{Intel oneAPI}套件，\textrm{BLAS}、\textrm{LAPACK}和\textrm{ScaLAPACK}和\textrm{FFTW}库均已包含在\textrm{MKL}库中，无须单独安装。如果使用\textrm{NVIDIA HPC-SDK}套装，\textrm{BLAS}、\textrm{LAPACK}和\textrm{ScaLAPACK}同样随\textrm{SDK}一同提供，唯一需要自行安装的数值库是\textrm{FFTW}。而开源的数学库可以在\url{https://www.netlib.org/liblist.html}找到全部各类库安装包。

%\section{\rm{MPI}库编译要求}
\textrm{VASP}支持分布式内存并行计算，通过\textrm{MPI}实现并行计算的信息同步与通信，\textrm{VASP}推荐的\textrm{MPI}库包括:
\begin{itemize}
	\item \textrm{Intel oneAPI MPI}(与\textrm{Intel}编译器配套)
	\item \textrm{OpenMPI}(最通用的开源\textrm{MPI}编译器和支持库，适用性较广)
	\item \textrm{NVIDIA HPC-SDK}自带的\textrm{OpenMPI}(面向\textrm{CUDA-aware}，可支持\textrm{GPU}间通信)
\end{itemize}
在高性能计算集群环境下，\textrm{MPI}属于标配，用户只需向集群运维和管理员确认编译器环境加载方案，遵照运行即可。用户可以通过类似\texttt{which mpiicc}等命令，检查相关\textrm{MPI}软件是否正确安装或加载。
%### 1.4 可选依赖库
%\section{可选的依赖库}

除去前述必要的组件外，\textrm{VASP}还支持若干可选功能，这些功能也需要额外的库支持。但因为这些功能用得较少，此处只对这些库作简单介绍。
\begin{itemize}
	\item \textrm{HDF5}:~用于支持读写\textrm{HDF5}格式文件(如\textrm{vaspout.h5})。%，**强烈推荐启用**，是py4vasp等后处理工具的前提
	\item \textrm{NCCL~(NVIDIA Collective Communications Library)}:~用于支持\textrm{GPU}多卡通信，%**强烈推荐**
		用于多\textrm{GPU}性能优化。
	\item \textrm{Wannier90}:~用于\textrm{Wannier}函数计算。
	\item \textrm{DFTD4}：用于\textrm{DFT-D4}色散校正。
\end{itemize}

\section{常见\rm{makefile.include}文件模版}
\textrm{VASP}源代码的\textrm{arch/}目录下，提供了大量的模板文件，根据文件名后缀，用户大致可以判断所使用的编译环境的组合形式:
\begin{itemize}
	\item \textrm{makefile.include.oneapi}:~ 适用于\textrm{Intel oneAPI}编译器 + \textrm{MKL}
	\item \textrm{makefile.include.intel}:~ 适用于早期版本\textrm{Intel Parallel Studio}
	\item \textrm{makefile.include.linux\_intel}:~ 适用于\textrm{Intel}编译器 + \textrm{MKL}
	\item \textrm{makefile.include.nvhpc\_omp\_acc}:~适用于\textrm{NVIDIA HPC-SDK} + \textrm{OpenACC} + \textrm{OpenMP}
	\item \textrm{makefile.include.linux\_nv\_acc}:~适用于\textrm{NVIDIA HPC-SDK} + \textrm{OpenACC}
	\item \textrm{makefile.include.aocc\_ompi\_aocl}:~适用于\textrm{AMD AOCC} + \textrm{OpenMPI} + \textrm{AOCL}
	\item \textrm{makefile.include.gfortran}:~适用于\textrm{GCC/GFortran}(最通用的基础版)
\end{itemize}
%**示例1：Intel oneAPI + MKL（纯CPU）** 
一般最常用的\textrm{makefile.include}模版文件有：
\begin{itemize}
\item 纯\textrm{CPU}架构:~\textrm{Intel oneAPI + MKL}
{\fontsize{8.2pt}{6.2pt}\selectfont{
	\lstinputlisting[linerange=1-56]{Figures/makefile.include.linux_intel_omp} %为保险:~选用文件名绝对路径
}}
	\item \textrm{GPU}架构:~\textrm{NVIDIA HPC-SDK + OpenAmpi}
{\fontsize{8.2pt}{6.2pt}\selectfont{
	\lstinputlisting[linerange=1-80]{Figures/makefile.include.linux_nv_omp} %为保险:~选用文件名绝对路径
}}
\item 纯开源编译:~\textrm{Linux}默认开源编译环境
{\fontsize{8.2pt}{6.2pt}\selectfont{
	\lstinputlisting[linerange=1-48]{Figures/makefile.include.linux_gfortran} %为保险:~选用文件名绝对路径
}}
\end{itemize}
%
在编译\textrm{VASP}软件前，用户从\textrm{arch/}目录复制最接近当前硬件配置的模版形式：
%#!/bin/bash
%echo "Hello from Bash"
%ls -la /tmp
%		\begin{lstlisting}
%   ```bash
%   ```
%\end{lstlisting}
%Welcome to LaTeX command line!
%
%\verb|latex document.tex|   # 编译 LaTeX 文档
%\verb|pdflatex document.tex|   # 编译为 PDF 文件
%
%\cmd{ls -l}   # 列出当前目录下的文件和目录
%\cmd{cd /path/to/directory}   # 进入指定目录
%\cmd{mkdir new_directory}   # 创建新目录
\begin{verbatim}
cp arch/makefile.include.your_choice ./makefile.include
\end{verbatim}

\section{\rm{makefile.include}的内容简介}
\textrm{makefile.include}是\textrm{VASP}编译中的核心配置文件，它指定了\textrm{make}命令%\footnote{\textrm{make}是\textrm{Linux}操作系统下的一个构建工具，通过执行\textrm{make}命令，解释\textrm{makefile}文件中指定的指令，完成自动化编译和构建软件项目。即\textrm{make}可用于开发人员管理和维护复杂的软件项目，自动化编译过程并生成可执行文件。}
编译\textrm{VASP}源代码所使用的各类编译器版本、位置、编译选项、所需的库函数，以及编译过程中各种库函数的链接关系等。%VASP官方明确表示：“从头编写makefile.include文件并不容易，因此建议从与系统最接近的模板文件开始”。

%### 6.1 模板文件的位置与选择
%**重要原则**：
%> **务必
习惯上，建议用户使用与所编译\textrm{VASP}版本一同发布的makefile.include模板文件，这些发布的编译安装文件中的环境变量和参数都是经过开发团队测试过的，旧版本的模板文件可能不完全适用新版本，反之亦然。%**
%

%### 6.2 makefile.include的结构
一个完整的\textrm{makefile.include}通常包含以下部分：
\begin{itemize}
	\item 预编译选项\textrm{(CPP\_OPTIONS)}
{\fontsize{8.2pt}{6.2pt}\selectfont{
	\lstinputlisting[linerange=1-12]{Figures/makefile.include.linux_intel_omp} %为保险:~选用文件名绝对路径
}}

%CPP_OPTIONS = -DHOST=\"LinuxIFC\" \
%              -DMPI -DMPI_BLOCK=8000 -Duse_collective \
%              -DscaLAPACK \
%              -DCACHE_SIZE=4000 \
%              -Davoidalloc \
%              -Dvasp6
%```
%

\begin{table}[h!]
    \centering
    \begin{tabular*}{0.70\textwidth}{|@{\extracolsep{\fill}}c|@{\extracolsep{\fill}}l|}
        \toprule
选项 &功能\\
\midrule
\textrm{-DMPI} &采用\textrm{MPI}并行\\
\textrm{-DMPI\_BLOCK=N} &指定\textrm{MPI}通信块大小\\
\textrm{-Duse\_collective} &指定使用集合通信\\
\textrm{-DscaLAPACK}  & 使用\textrm{ScaLAPACK}库\\
\textrm{-DCACHE\_SIZE=N} &设定缓存大小\\~%(限于内置\textrm{FFT}例程)
\textrm{-Davoidalloc} &该选项为避免出现某些内存分配问题\\
\textrm{-Dvasp6} &\textrm{VASP 6.x}~版本标识\\
\textrm{-DVASP\_HDF5} & 使用可选的\textrm{HDF5}库支持\\
\textrm{-D\_OPENACC} & 指定\textrm{GPU}版本的\textrm{OpenACC}\\
\textrm{-DUSENCCL} & 支持多\textrm{GPU}通信的\textrm{NCCL}\\
        \bottomrule
    \end{tabular*}
    \caption{主要预编译选项说明}
\label{Tab:Makefile-Options}
\end{table}

%#### （二）
\item 编译器及编译选项
{\fontsize{8.2pt}{6.2pt}\selectfont{
	\lstinputlisting[linerange=13-24]{Figures/makefile.include.linux_intel_omp} %为保险:~选用文件名绝对路径
}}

%```makefile
%CPP         = nvfortran -Mpreprocess -Mfree -Mextend -E $(CPP_OPTIONS) $*$(FUFFIX) > $*$(SUFFIX)
%FC          = mpif90 -acc -gpu=cc60,cc70,cc80,cuda11.0 -mp
%FCL         = mpif90 -acc -gpu=cc60,cc70,cc80,cuda11.0 -mp -c++libs
%FREE        = -Mfree
%FFLAGS      = -Mbackslash -Mlarge_arrays
%OFLAG       = -fast
%DEBUG       = -Mfree -O0 -traceback
%```

\begin{table}[h!]
    \centering
    \begin{tabular*}{0.55\textwidth}{|@{\extracolsep{\fill}}c|@{\extracolsep{\fill}}l|}
        \toprule
选项 &功能\\
\midrule
\textrm{CPP} &预处理命令\\
\textrm{FC} &指定\textrm{Fortran}代码编译器\\
\textrm{FCL} &用于\textrm{Fortran}代码链接指令\\
\textrm{FREE}  &格式自由的源代码编译选项\\
\textrm{FFLAGS} &指定的通用编译选项\\
\textrm{OFLAG} &用于代码优化的选项\\
\textrm{DEBUG} &用于代码调试的选项\\
        \bottomrule
    \end{tabular*}
    \caption{主要编译器和编译选项的说明}
\label{Tab:Makefile-Options}
\end{table}


%**优化级别选择**：
%| 级别 | 选项 | 适用场景 |
%|------|------|---------|
%| 调试 | `-O0` | 开发调试 |
%| 基础 | `-O2` | 通用计算 |
%| 高性能 | `-O3` | 生产计算 |
%| 极致 | `-fast`（NVHPC）或`-xHOST -O3`（Intel） | 追求极致性能 |
%
%#### （三）链接库配置
\item 库函数链接配置
{\fontsize{8.2pt}{6.2pt}\selectfont{
	\lstinputlisting[linerange=25-39]{Figures/makefile.include.linux_intel_omp} %为保险:~选用文件名绝对路径
}}

%```makefile
%OBJECTS     = fftmpiw.o fftmpi_map.o fftw3d.o fft3dlib.o
%LLIBS       = -cudalib=cublas,cusolver,cufft,nccl -cuda
%```
\begin{table}[h!]
    \centering
    \begin{tabular*}{0.95\textwidth}{|@{\extracolsep{\fill}}c|@{\extracolsep{\fill}}l|}
        \toprule
选项 &功能\\
\midrule
\textrm{MKL\_LIBS} &指定重要数学库的存储路径\\
\textrm{BLAS/BLACS} &\makecell[tl]{涉及向量运算、矩阵乘法的基本线性代数运算库\\ (\textrm{BLAVS}支持并行计算)}\\
\textrm{LAPACK/SCALAPACK} &\makecell[tl]{涉及矩阵对角化和本征值求解的线性代数运算库\\ (\textrm{SCALPACK}支持并行计算)}\\
\textrm{OBJECTS}  &支持\textrm{FFTW}所需要的各类数学库的目标文件和库函数\\
\textrm{INCS} &\textrm{FFTW}库中必要的头文件目录\\
\textrm{OBJECTS} &用于代码优化的选项\\
\textrm{LLIBS} &指定了所需要的全部库函数\\
%后续编译过程中，链接器将通过\textrm{\$(FCL) -o vasp ..all-objects.. \$(LLIBS)}命令完成最终全部库函数的链接
        \bottomrule
    \end{tabular*}
    \caption{主要链接的库函数说明}
\label{Tab:Makefile-Options}
\end{table}

%**优化级别选择**：

%`LLIBS`指定了需要链接的所有库。
\end{itemize}

用户编译和安装\textrm{VASP}软件，主要就是根据硬件配置，设置和修改\textrm{makefile.include}文件中的一些参数和变量。具体操作步骤依次为：
\begin{enumerate}
	\item  指定编译器的正确路径，如:\\
{\fontsize{8.2pt}{6.2pt}\selectfont{
	\lstinputlisting[linerange=15-16]{Figures/makefile.include.linux_intel_omp} %为保险:~选用文件名绝对路径
}}
此处示例的\textrm{FC}和\textrm{FCL}指向的编译器为\textrm{mpiifort}，它应该是已加载过的环境变量(通过命令\texttt{which mpiifort}可以获取)。用户能使用的是当前操作系统安装和配置的编译器，一般可以通过\texttt{which mpif90}或\texttt{which ifort}等命令检查，如果指定的编译器已经安装，但没有加载在环境变量内，则必须设置编译器的绝对路径。
	\item 设置各类库函数路径:~包括\textrm{MKL\_ROOT}、\textrm{FFTW\_ROOT}、\textrm{NVROOT}等并指定必要的库函数名，如:\\
{\fontsize{8.2pt}{6.2pt}\selectfont{
	\lstinputlisting[linerange=25-39]{Figures/makefile.include.linux_intel_omp} %为保险:~选用文件名绝对路径
}}
此处用户设定\textrm{MKL\_PATH}的路径来设置\textrm{MKL}库的路径和变量，原则上是在\textrm{Intel oneAPI}套装完整安装并加载后使用，即通过命令\texttt{echo \$MKLROOT}可以获取变量\textrm{\$MKLROOT}的具体路径，此外，\textrm{\$MKROOT}也包含了\textrm{FFTW}的相关库。后缀名为\textrm{.o}的文件，表示编译后链接的目标对象，而后缀名\textrm{.a}的表示静态库文件，一般要求写出完整路径；而以\textrm{-l}引导的\textrm{-lmkl\_xxxx}，表示使用\textrm{\$MKL\_PATH}目录下的\textrm{libmkl\_xxxx.so}库(后缀名为\textrm{.so}的表示动态库文件)。如果用户使用手动安装相关数学库，也要遵照上述规则，指定相应的库函数的具体名称和位置。%比较老的\textrm{VASP}版本
%   - 确认`FFTW_ROOT`或`FFTW`路径正确
%   - 确认`NVROOT`正确（NVIDIA HPC-SDK）
	\item 选择预编译选项:~根据需要启用或禁用功能:
{\fontsize{8.2pt}{6.2pt}\selectfont{
	\lstinputlisting[linerange=40-56]{Figures/makefile.include.linux_intel_omp} %为保险:~选用文件名绝对路径
}}
一般地，此处及以下，用户不作修改。但用户可以根据这里的参数设置了解预编译及选项参数的设置方案。
	\item 调整优化级别:~根据硬件和需求调整\textrm{OFLAG}。\\
		如果需要对\textrm{VASP}软件的代码进行进一步测试或开发，可以修改相关编译代码的优化层级。需要指出的是，\textrm{VASP}软件开发团队对于\textrm{VASP}代码的修改权限，做出严格的限制。
\end{enumerate}

原则上，如果用户的环境变量已经正确加载，只要在\textrm{makefile.include}文件中修改上述几处内容，即可完成\textrm{VASP}软件的正确编译。编译命令为\texttt{make all}。根据硬件的不同，一般编译时间为20\textrm{min}-45\textrm{min}不等，编译后产生的\textrm{VASP}可执行文件位于目录\textrm{bin/}下，版本依次为：
\begin{itemize}
	\item \textrm{vasp\_std}:~标准版，最通用的\textrm{VASP}版本。
	\item \textrm{vasp\_gam}:~允许仅有单个$\vec k$点(默认为$\Gamma$点)的\textrm{VASP}计算。
	\item \textrm{vasp\_ncl}:~实现非线性光学计算的\textrm{VASP}计算。
\end{itemize}
将可执行文件所在目录\textrm{bin/}，加载到环境变量并生效后，即可在输入文件所在目录下运行可执行文件，直接启动\textrm{VASP}计算。在高性能计算环境下，可在计算目录(输入文件所在目录)的作业提交脚本中，直接指定\textrm{VASP}可执行文件路径后，提交计算任务。
%```makefile
%CPP_OPTIONS = -DHOST=\"LinuxIFC\" -DMPI -DMPI_BLOCK=8000 -Duse_collective -DscaLAPACK -DCACHE_SIZE=4000 -Davoidalloc -Dvasp6
%CPP         = mpiifort -E -DMPI -C -P $(CPP_OPTIONS) $< >$*.F90
%FC          = mpiifort -I$(MKLROOT)/include
%FCL         = mpiifort -L$(MKLROOT)/lib
%FREE        = -free -names lowercase
%FFLAGS      = -assume byterecl -O2 -fp-model precise
%OFLAG       = -O3 -xHOST
%OBJECTS     = fftmpiw.o fftmpi_map.o fftw3d.o fft3dlib.o
%LLIBS       = -lmkl_scalapack_lp64 -lmkl_blacs_intelmpi_lp64 -lmkl_intel_lp64 -lmkl_sequential -lmkl_core
%```
%
%**示例2：NVIDIA HPC-SDK + OpenACC（GPU）** 
%
%```makefile
%CPP_OPTIONS = -DHOST=\"LinuxNV\" -DMPI -DMPI_BLOCK=8000 -Duse_collective -DscaLAPACK -DCACHE_SIZE=4000 -Davoidalloc -Dvasp6 -D_OPENMP -D_OPENACC -DUSENCCL -DUSENCCLP2P
%CPP         = nvfortran -Mpreprocess -Mfree -Mextend -E $(CPP_OPTIONS) $*$(FUFFIX) > $*$(SUFFIX)
%FC          = mpif90 -acc -gpu=cc60,cc70,cc80,cuda11.0 -mp
%FCL         = mpif90 -acc -gpu=cc60,cc70,cc80,cuda11.0 -mp -c++libs
%FREE        = -Mfree
%FFLAGS      = -Mbackslash -Mlarge_arrays
%OFLAG       = -fast
%OBJECTS     = fftmpiw.o fftmpi_map.o fftw3d.o fft3dlib.o
%LLIBS       = -cudalib=cublas,cusolver,cufft,nccl -cuda
%QD          = $(NVROOT)/compilers/extras/qd
%```
%
%%## 二、编译器详解
%
%编译器是VASP编译中最核心的决策。不同的编译器在优化能力、兼容性和对特定硬件（Intel/AMD/NVIDIA）的支持上各有优劣。以下逐一详解各主流编译器的特点与配置方法。
%
%### 2.1 Intel oneAPI编译器
%
%Intel oneAPI是VASP编译的**最经典、最成熟**的选择，多年来一直是VASP官方推荐的首选编译器套件。
%
%**核心组件**：
%- **ifort**（传统Fortran编译器）或**ifx**（基于LLVM的新一代Fortran编译器）
%- **icc/icx**（C/C++编译器）
%- **MKL**（Math Kernel Library，包含BLAS、LAPACK、ScaLAPACK及FFTW接口）
%
%**优势**：
%- 与VASP代码的兼容性经过长期验证
%- MKL数学库在Intel CPU上的性能优化极为出色
%- 提供了完整的工具链，无需额外安装BLAS/LAPACK/ScaLAPACK
%
%**典型makefile.include配置**（VASP 6.5及以上版本使用`makefile.include.oneapi`）：
%
%```makefile
%# Default precompiler options
%CPP_OPTIONS = -DHOST=\"LinuxIFC\" \
%              -DMPI -DMPI_BLOCK=8000 -Duse_collective \
%              -DscaLAPACK \
%              -DCACHE_SIZE=4000 \
%              -Davoidalloc \
%              -Dvasp6 \
%              -Duse_bse_te \
%              -Dtbdyn \
%              -Dqd_emulate \
%              -Dfock_dblbuf
%
%CPP         = mpiifort -E -DMPI -C -P $(CPP_OPTIONS) $< >$*.F90
%FC          = mpiifort -I$(MKLROOT)/include
%FCL         = mpiifort -L$(MKLROOT)/lib
%
%FREE        = -free -names lowercase
%FFLAGS      = -assume byterecl -O2 -fp-model precise
%OFLAG       = -O3 -xHOST
%
%MKL_PATH    = $(MKLROOT)/lib/intel64
%BLAS        = -L$(MKL_PATH) -lmkl_intel_lp64 -lmkl_sequential -lmkl_core
%LAPACK      = $(BLAS)
%SCALAPACK   = $(MKL_PATH)/libmkl_scalapack_lp64.a -Wl,--start-group \
%              $(MKL_PATH)/libmkl_intel_lp64.a \
%              $(MKL_PATH)/libmkl_sequential.a \
%              $(MKL_PATH)/libmkl_core.a \
%              $(MKL_PATH)/libmkl_blacs_intelmpi_lp64.a \
%              -Wl,--end-group
%```
%
%**关键编译选项说明**：
%- `-xHOST`：自动检测当前CPU架构并生成最优指令集
%- `-O3`：最高级别优化
%- `-fp-model precise`：保证浮点运算的可重现性
%- `-assume byterecl`：解决记录长度相关的兼容性问题
%- `-mkl=sequential`：使用MKL的串行版本（若需要多线程则改为`-mkl=parallel`）
%
%### 2.2 NVIDIA HPC-SDK编译器
%
%NVIDIA HPC-SDK是**GPU加速版VASP编译的唯一选择**。自VASP 6.2.0正式发布OpenACC GPU移植版以来，NVIDIA HPC-SDK已成为GPU计算场景下的标准工具链。
%
%**核心组件**：
%- **nvfortran**（Fortran编译器，支持OpenACC）
%- **nvc/nvc++**（C/C++编译器）
%- **CUDA Toolkit**（包含CUDA运行时、cuBLAS、cuSOLVER等）
%- **NCCL**（多GPU通信库）
%- **OpenMPI**（CUDA-aware MPI实现）
%
%**版本要求**：
%- NVIDIA HPC-SDK版本必须 **≥ 21.2**
%- **特别注意**：版本22.1和22.2存在严重bug，无法在启用OpenMP线程的情况下执行OpenACC版本，应避免使用
%- 官方建议使用20.9版本
%
%**GPU硬件要求**：
%VASP官方已测试的GPU包括：
%- **数据中心GPU**：P100（Pascal）、V100（Volta）、A100（Ampere）
%- **工作站GPU**：GP100（Pascal）、GV100（Volta）
%- **游戏卡**：技术上可行但**不推荐**——双精度浮点（FP64）性能不足，且显存通常不带ECC纠错
%
%**典型makefile.include配置**（`makefile.include.nvhpc_omp_acc`）：
%
%```makefile
%# Default precompiler options
%CPP_OPTIONS = -DHOST=\"LinuxNV\" \
%              -DMPI -DMPI_BLOCK=8000 -Duse_collective \
%              -DscaLAPACK \
%              -DCACHE_SIZE=4000 \
%              -Davoidalloc \
%              -Dvasp6 \
%              -Duse_bse_te \
%              -Dtbdyn \
%              -Dqd_emulate \
%              -Dfock_dblbuf \
%              -D_OPENMP \
%              -D_OPENACC \
%              -DUSENCCL -DUSENCCLP2P
%
%CPP         = nvfortran -Mpreprocess -Mfree -Mextend -E $(CPP_OPTIONS) $*$(FUFFIX) > $*$(SUFFIX)
%
%# N.B.: you might need to change the cuda-version here
%# to one that comes with your NVIDIA-HPC SDK
%FC          = mpif90 -acc -gpu=cc60,cc70,cc80,cuda11.0 -mp
%FCL         = mpif90 -acc -gpu=cc60,cc70,cc80,cuda11.0 -mp -c++libs
%
%FREE        = -Mfree
%FFLAGS      = -Mbackslash -Mlarge_arrays
%OFLAG       = -fast
%
%OBJECTS     = fftmpiw.o fftmpi_map.o fftw3d.o fft3dlib.o
%LLIBS       = -cudalib=cublas,cusolver,cufft,nccl -cuda
%
%# QD (quadruple precision emulation)
%QD          = $(NVROOT)/compilers/extras/qd
%```
%
%**关键编译选项说明**：
%- `-acc`：启用OpenACC编译
%- `-gpu=cc60,cc70,cc80,cuda11.0`：指定GPU计算能力和CUDA版本
%- `-mp`：启用OpenMP多线程支持
%- `-c++libs`：链接C++标准库
%- `-fast`：相当于`-O3 -Mipa=fast -Mvect -Mfprelaxed -Mautoinline`等优化选项的组合
%- `-cudalib=cublas,cusolver,cufft,nccl`：链接CUDA数学库和NCCL
%
%### 2.3 GCC/GFortran编译器
%
%GCC是Linux系统的默认编译器套件，对于预算有限或无法获取商业编译器的用户来说，是一个可行的选择。
%
%**版本要求**：
%- 仅从**VASP 6.2.0**开始支持GCC 7.X以上版本
%- 早期VASP版本中的某些代码结构与GCC的严格检查不兼容
%
%**优势**：
%- 完全开源，无需商业许可
%- 与各种Linux发行版兼容性好
%- 适合个人电脑或小型集群
%
%**劣势**：
%- 优化能力不如Intel编译器
%- 数学库需要单独安装（OpenBLAS或ATLAS）
%- 对AMD CPU的支持不如AOCC
%
%### 2.4 AMD AOCC编译器
%
%AOCC（AMD Optimizing C/C++ and Fortran Compiler）是AMD专门为其EPYC处理器优化的编译器套件。
%
%**适用场景**：
%- 搭载AMD EPYC处理器的计算集群
%- 追求AMD平台上的极致性能
%
%**配套数学库**：
%- **AOCL**（AMD Optimizing CPU Libraries）：包含AMD优化的BLAS（amdblis）、LAPACK（amdlapack）和FFTW（amdfftw）实现
%
%**典型配置**：
%VASP官方提供了`makefile.include.aocc_ompi_aocl`模板，用于AOCC + OpenMPI + AOCL的组合编译。
%
%### 2.5 编译器选择决策指南
%
%| 场景 | 推荐编译器 | 理由 |
%|------|-----------|------|
%| Intel CPU + 纯CPU计算 | **Intel oneAPI** | MKL性能最优，兼容性最好 |
%| NVIDIA GPU计算 | **NVIDIA HPC-SDK** | 唯一支持OpenACC GPU编译的选择 |
%| AMD EPYC CPU | **AMD AOCC** | 针对AMD架构深度优化 |
%| 个人电脑/学习用途 | **GCC/GFortran** | 开源免费，易于获取 |
%| 混合架构（Intel+GPU） | **NVIDIA HPC-SDK + MKL** | 可用MKL作为数学库，HPC-SDK负责GPU编译 |
%
%
%## 三、关键数学库详解
%
%VASP的数学库是其计算性能的基石。以下逐一详解各数学库的作用、安装与配置方法。
%
%### 3.1 BLAS（Basic Linear Algebra Subprograms）
%
%**功能**：BLAS是基础线性代数运算的标准接口，提供向量-向量运算（Level 1）、矩阵-向量运算（Level 2）和矩阵-矩阵运算（Level 3）。VASP的大量计算（如哈密顿量构建、波函数更新等）都依赖BLAS。
%
%**可选实现**：
%| 实现 | 特点 | 适用场景 |
%|------|------|---------|
%| **Intel MKL** | 性能最优，Intel CPU上自动使用AVX-512 | Intel CPU |
%| **OpenBLAS** | 开源，性能优秀，支持多种架构 | 通用 |
%| **AMD BLIS** | AMD优化，AOCL的一部分 | AMD EPYC |
%| **Netlib BLAS** | 参考实现，性能较低 | 仅用于测试 |
%
%**在makefile.include中的配置**（以Intel MKL为例）：
%```makefile
%BLAS = -L$(MKL_PATH) -lmkl_intel_lp64 -lmkl_sequential -lmkl_core
%```
%或简化为：
%```makefile
%FCL = $(FC) -mkl=sequential
%```
%
%### 3.2 LAPACK（Linear Algebra PACKage）
%
%**功能**：LAPACK在BLAS之上提供了更高级的线性代数功能，包括求解线性方程组、特征值分解、奇异值分解（SVD）、QR分解等。VASP在对角化哈密顿量时重度依赖LAPACK。
%
%**与BLAS的关系**：LAPACK本身调用BLAS完成底层运算，因此LAPACK的性能高度依赖于底层BLAS的优化程度。
%
%**在makefile.include中的配置**：
%大多数情况下，LAPACK与BLAS使用相同的库路径。在Intel MKL中：
%```makefile
%LAPACK = $(BLAS)
%```
%
%### 3.3 ScaLAPACK（Scalable LAPACK）
%
%**功能**：ScaLAPACK是LAPACK的分布式内存并行版本，支持将大规模矩阵运算分布在多个MPI进程上并行执行。对于大型体系（如数百个原子以上的计算），ScaLAPACK是提升并行效率的关键。
%
%**依赖关系**：ScaLAPACK依赖于BLACS（Basic Linear Algebra Communication Subprograms）提供通信功能。
%
%**在makefile.include中的配置**：
%- **Intel MKL**（包含ScaLAPACK和BLACS）：
%```makefile
%SCALAPACK = -lmkl_scalapack_lp64.a -lmkl_blacs_openmpi_lp64
%```
%- **独立安装的ScaLAPACK**：
%```makefile
%SCALAPACK_ROOT = /path/to/scalapack
%SCALAPACK = -L$(SCALAPACK_ROOT)/lib -lscalapack
%```
%
%**启用ScaLAPACK的预编译选项**：
%```makefile
%CPP_OPTIONS += -DscaLAPACK
%```
%
%### 3.4 FFTW（Fastest Fourier Transform in the West）
%
%**功能**：FFTW是实现快速傅里叶变换（FFT）的业界标准库。在VASP中，FFT用于在实空间和倒空间之间转换波函数和电荷密度——这是平面波赝势方法中最核心、最频繁的操作之一。
%
%**为什么FFTW是必需的**：
%VASP基于平面波基组，波函数在倒空间中表示，而交换关联势、电荷密度等物理量在实空间中计算更为方便。每次迭代都需要在实空间和倒空间之间反复转换，FFT的效率直接决定了整个计算的性能。
%
%**重要说明**：
%- 如果使用**NVIDIA HPC-SDK**，FFTW是**唯一需要自行安装**的数值库（BLAS、LAPACK、ScaLAPACK已包含在SDK中）
%- 如果使用**Intel MKL**，MKL提供了FFTW的接口封装（`-mkl`选项会自动链接MKL的FFT实现）
%
%**FFTW的安装**：
%```bash
%# 下载FFTW源代码
%wget http://www.fftw.org/fftw-3.3.10.tar.gz
%tar -xzf fftw-3.3.10.tar.gz
%cd fftw-3.3.10
%
%# 配置（以Intel编译器为例）
%./configure --prefix=/path/to/fftw \
%            --enable-mpi \
%            --enable-openmp \
%            --enable-threads \
%            --enable-shared \
%            CC=icc FC=ifort F77=ifort
%
%make -j8
%make install
%```
%
%**关键配置选项说明**：
%- `--enable-mpi`：启用MPI并行FFT
%- `--enable-openmp`：启用OpenMP多线程FFT
%- `--enable-threads`：启用POSIX线程支持
%- `--enable-shared`：生成共享库（便于动态链接）
%
%**在makefile.include中的配置**：
%```makefile
%FFTW_ROOT = /path/to/fftw
%LLIBS += -L$(FFTW_ROOT)/lib -lfftw3 -lfftw3_mpi -lfftw3_omp
%INCS += -I$(FFTW_ROOT)/include
%```
%
%**使用MKL的FFTW接口**（更简便）：
%```makefile
%FCL = $(FC) -mkl=sequential
%OBJECTS = fftmpiw.o fftmpi_map.o fftw3d.o fft3dlib.o
%```
%
%### 3.5 数学库配置的典型组合
%
%| 编译器 | 数学库组合 | 配置方式 |
%|--------|-----------|---------|
%| Intel oneAPI | MKL（含BLAS+LAPACK+ScaLAPACK+FFTW） | `-mkl=sequential` |
%| NVIDIA HPC-SDK | SDK内置BLAS/LAPACK/ScaLAPACK + 自行安装FFTW | 分别指定路径 |
%| GCC + OpenBLAS | OpenBLAS + LAPACK + 独立ScaLAPACK + FFTW | 逐个指定 |
%| AMD AOCC | AOCL（含BLIS+LAPACK+amdfftw） | 使用AOCL路径 |
%
%
%## 四、MPI库详解
%
%MPI（Message Passing Interface）是VASP实现分布式内存并行（多节点/多进程）计算的核心通信库。
%
%### 4.1 MPI库的选型
%
%| MPI实现 | 特点 | 适用场景 |
%|---------|------|---------|
%| **Intel MPI** | 与Intel编译器深度集成，性能优异 | Intel CPU集群 |
%| **OpenMPI** | 开源，适用性最广，支持CUDA-aware | 通用，GPU集群 |
%| **NVIDIA HPC-SDK自带OpenMPI** | CUDA-aware，与GPU编译无缝集成 | NVIDIA GPU集群 |
%
%### 4.2 MPI与编译器的配合
%
%MPI库必须与Fortran编译器**配套使用**——用哪个编译器编译的MPI库，就必须用哪个编译器来编译VASP。混用不同编译器编译的MPI库会导致链接错误或运行时崩溃。
%
%**编译器与MPI的对应关系**：
%- Intel编译器 → Intel MPI 或 OpenMPI（用Intel编译器编译）
%- NVIDIA HPC-SDK → 自带OpenMPI（CUDA-aware）
%- GCC → OpenMPI（用GCC编译）
%
%### 4.3 MPI在makefile.include中的配置
%
%**使用Intel MPI**：
%```makefile
%FC = mpiifort
%FCL = mpiifort
%```
%
%**使用OpenMPI**：
%```makefile
%FC = mpif90
%FCL = mpif90
%```
%
%**使用NVIDIA HPC-SDK自带的OpenMPI**：
%```makefile
%FC = mpif90 -acc -gpu=cc60,cc70,cc80,cuda11.0 -mp
%FCL = mpif90 -acc -gpu=cc60,cc70,cc80,cuda11.0 -mp -c++libs
%```
%
%### 4.4 MPI相关的预编译选项
%
%在makefile.include的`CPP_OPTIONS`中，以下选项控制MPI相关功能：
%
%```makefile
%-DMPI              # 启用MPI支持（必需）
%-DMPI_BLOCK=8000   # MPI通信块大小
%-Duse_collective   # 使用集合通信（优化并行效率）
%```
%
%### 4.5 CUDA-aware MPI
%
%对于GPU计算，MPI库必须是**CUDA-aware**的，即能够直接发送和接收GPU显存中的数据，而无需先将数据拷贝到CPU内存。
%
%NVIDIA HPC-SDK自带的OpenMPI是CUDA-aware的，这也是官方推荐在GPU场景下使用NVIDIA HPC-SDK的重要原因之一。
%
%
%## 五、GPU编译详解
%
%自VASP 6.2.0起，VASP官方正式发布了基于**OpenACC**标准的GPU加速版本。在此之前存在的CUDA-C GPU移植版已被**废弃**，自VASP 6.3.0起已完全移除。因此，**当前所有GPU编译都应基于OpenACC路线**。
%
%### 5.1 OpenACC vs CUDA-C
%
%| 特性 | OpenACC版本 | CUDA-C版本（已废弃） |
%|------|------------|---------------------|
%| 状态 | **官方推荐，持续开发** | VASP 6.3.0起已移除 |
%| 编译器 | NVIDIA HPC-SDK（≥21.2） | 旧版CUDA Toolkit |
%| 编程模型 | 指令式（编译器自动并行） | 显式GPU编程 |
%| 可维护性 | 代码统一，易于维护 | 需要维护两套代码 |
%
%### 5.2 GPU编译的软件栈
%
%完整的GPU编译需要以下组件：
%
%1. **NVIDIA HPC-SDK（≥21.2）** ：提供nvfortran编译器（支持OpenACC）
%2. **CUDA Toolkit（≥10.0）** ：包含CUDA驱动和运行时——已包含在HPC-SDK中
%3. **CUDA-aware MPI**：已包含在HPC-SDK的OpenMPI中
%4. **NCCL（≥2.7.8）** ：多GPU通信库——**强烈推荐**，已包含在HPC-SDK中
%5. **FFTW**：**唯一需要自行安装的库**
%
%### 5.3 OpenACC编译的关键选项
%
%在NVIDIA HPC-SDK中，启用GPU编译的核心选项是：
%
%```makefile
%FC = mpif90 -acc -gpu=cc60,cc70,cc80,cuda11.0 -mp
%```
%
%**各选项含义**：
%| 选项 | 含义 |
%|------|------|
%| `-acc` | 启用OpenACC编译 |
%| `-gpu=cc60,cc70,cc80,cuda11.0` | 指定GPU计算能力和CUDA版本 |
%| `-mp` | 启用OpenMP（用于CPU多线程） |
%
%**GPU计算能力（Compute Capability）对照**：
%| 架构 | 计算能力 | -gpu参数 |
%|------|---------|----------|
%| Pascal | 6.0/6.1 | `cc60` |
%| Volta | 7.0 | `cc70` |
%| Turing | 7.5 | `cc75` |
%| Ampere | 8.0 | `cc80` |
%| Ada Lovelace | 8.9 | `cc89` |
%| Hopper | 9.0 | `cc90` |
%
%### 5.4 NCCL与多GPU通信
%
%NCCL（NVIDIA Collective Communications Library）是实现多GPU之间高效通信的关键库。在makefile.include中启用NCCL：
%
%```makefile
%CPP_OPTIONS += -DUSENCCL -DUSENCCLP2P
%LLIBS += -cudalib=nccl
%```
%
%**重要限制**：由于NCCL库的限制，OpenACC GPU版本目前**只能单进程运行**（即每个GPU一个MPI进程）。作为补充，可以通过OpenMP实现单进程内的多线程并行。
%
%### 5.5 GPU编译的典型工作流程
%
%**Step 1：安装NVIDIA HPC-SDK**
%
%```bash
%# 下载并安装HPC-SDK（以20.9为例）
%wget https://developer.download.nvidia.com/hpc-sdk/20.9/nvhpc_2020_209_Linux_x86_64_cuda_11.0.tar.gz
%tar -xzf nvhpc_2020_209_Linux_x86_64_cuda_11.0.tar.gz
%./nvhpc_2020_209_Linux_x86_64_cuda_11.0/install
%# 按提示指定安装路径，如 /opt/nvidia/hpc_sdk
%```
%
%**Step 2：设置环境变量**
%
%```bash
%export NVARCH=`uname -s`_`uname -m`
%export NVCOMPILERS=/opt/nvidia/hpc_sdk
%export PATH=$NVCOMPILERS/$NVARCH/20.9/compilers/bin:$PATH
%export LD_LIBRARY_PATH=$NVCOMPILERS/$NVARCH/20.9/compilers/lib:$LD_LIBRARY_PATH
%export PATH=$NVCOMPILERS/$NVARCH/20.9/comm_libs/mpi/bin:$PATH
%```
%
%**Step 3：准备makefile.include**
%
%```bash
%cp arch/makefile.include.linux_nv_acc ./makefile.include
%# 或使用OpenMP+OpenACC混合版本
%cp arch/makefile.include.nvhpc_omp_acc ./makefile.include
%```
%
%**Step 4：修改makefile.include**
%- 确认编译器路径正确
%- 确认CUDA版本与HPC-SDK匹配
%- 确认FFTW路径正确
%- 根据GPU型号调整`-gpu`参数
%
%**Step 5：编译**
%
%```bash
%make DEPS=1 -jN all
%```
%
%**特别说明**：OpenACC版本编译后得到的`vasp_std`、`vasp_gam`、`vasp_ncl`**直接支持GPU计算**，无需单独的`vasp_gpu`可执行文件。
%
%### 5.6 GPU编译的注意事项
%
%1. **双精度性能**：VASP计算主要依赖双精度浮点运算（FP64）。消费级游戏卡的FP64性能通常被大幅削减，**不推荐用于生产计算**。
%
%2. **显存容量**：VASP的GPU计算对显存需求较高，2GB显存很快就会填满。建议至少使用16GB以上显存的GPU。
%
%3. **NCORE参数**：在OpenACC版本中，INCAR中的`NCORE`**只能设为1**。
%
%4. **避免版本22.1/22.2**：这两个版本的NVIDIA HPC-SDK存在严重bug。
%
%5. **环境隔离**：建议将GPU编译的环境变量与Intel编译器环境分开，避免`mpirun`冲突。
%
%
%
%## 七、VASP的编译与安装
%
%### 7.1 标准编译流程
%
%**Step 1：获取源代码**
%从VASP官方门户下载源代码包，解压到目标位置。
%
%**Step 2：准备makefile.include**
%```bash
%cd /path/to/vasp.x.x.x
%cp arch/makefile.include.your_choice ./makefile.include
%# 根据系统情况编辑makefile.include
%```
%
%**Step 3：执行编译**
%```bash
%make DEPS=1 -jN <target>
%```
%
%**编译目标（target）** ：
%| 目标 | 生成的可执行文件 | 用途 |
%|------|-----------------|------|
%| `std` | `vasp_std` | 标准版本（最常用） |
%| `gam` | `vasp_gam` | Gamma-only版本（仅适用于Gamma点计算，效率更高） |
%| `ncl` | `vasp_ncl` | 非共线磁性版本 |
%| `all` | 全部三个 | 一次性编译所有版本 |
%
%**Step 4：测试**
%```bash
%./bin/vasp_std
%```
%如果输出显示VASP的版本信息和正常启动信息，说明编译成功。
%
%### 7.2 编译优化建议
%
%1. **并行编译**：使用`-jN`参数（N为并行任务数）加速编译过程
%2. **分步编译**：先编译`std`测试，确认无误后再编译`all`
%3. **保留编译日志**：重定向输出以便排查错误
%   ```bash
%   make all >& make.log
%   ```
%
%### 7.3 安装后的文件结构
%
%编译完成后，可执行文件位于`/path/to/vasp.x.x.x/bin/`目录下：
%```
%bin/
%├── vasp_std
%├── vasp_gam
%└── vasp_ncl
%```
%
%
%## 八、常见编译问题与解决方案
%
%### 8.1 编译器相关问题
%
%**问题1：`ifort: command not found`**
%- **原因**：Intel编译器未安装或环境变量未设置
%- **解决**：安装Intel oneAPI并运行`source /opt/intel/oneapi/setvars.sh`
%
%**问题2：`mpif90: command not found`**
%- **原因**：MPI库未安装或未加载
%- **解决**：安装OpenMPI或Intel MPI，并加载相应模块
%
%**问题3：GCC版本不兼容**
%- **原因**：VASP 6.2.0以下版本不支持GCC 7.X以上
%- **解决**：使用较低版本GCC，或升级到VASP 6.2.0以上
%
%### 8.2 数学库相关问题
%
%**问题4：`cannot find -lmkl_intel_lp64`**
%- **原因**：MKL库路径不正确
%- **解决**：确认`MKLROOT`环境变量已设置，检查`MKL_PATH`是否正确
%
%**问题5：FFTW链接错误**
%- **原因**：FFTW未安装或路径不正确
%- **解决**：安装FFTW并确认`FFTW_ROOT`路径正确
%
%**问题6：ScaLAPACK未找到**
%- **原因**：ScaLAPACK库未正确链接
%- **解决**：确保`CPP_OPTIONS`包含`-DscaLAPACK`，且`SCALAPACK`变量正确
%
%### 8.3 GPU编译相关问题
%
%**问题7：`-acc: command not found`**
%- **原因**：未使用NVIDIA HPC-SDK的编译器
%- **解决**：确认`FC`使用的是`nvfortran`或`mpif90`（来自HPC-SDK）
%
%**问题8：GPU计算能力不匹配**
%- **原因**：`-gpu`参数指定的计算能力与GPU不匹配
%- **解决**：查询GPU的计算能力，修改`-gpu`参数
%
%**问题9：NCCL相关错误**
%- **原因**：NCCL未安装或版本过低
%- **解决**：确认NCCL已包含在HPC-SDK中，或单独安装NCCL（≥2.7.8）
%
%**问题10：OpenACC与OpenMP冲突（22.1/22.2版本）**
%- **原因**：NVIDIA HPC-SDK 22.1/22.2存在bug
%- **解决**：升级到22.3以上，或编译时不启用OpenMP
%
%### 8.4 MPI相关问题
%
%**问题11：MPI版本不匹配**
%- **原因**：编译VASP时使用的MPI与运行时加载的MPI不一致
%- **解决**：确保编译和运行时使用相同的MPI实现和版本
%
%**问题12：CUDA-aware MPI未启用**
%- **原因**：使用的MPI实现不支持CUDA-aware
%- **解决**：使用NVIDIA HPC-SDK自带的OpenMPI
%
%### 8.5 调试技巧
%
%1. **查看详细错误信息**：
%   ```bash
%   make clean
%   make all >& make.log 2>&1
%   tail -100 make.log
%   ```
%
%2. **检查库依赖**：
%   ```bash
%   ldd bin/vasp_std
%   ```
%
%3. **验证GPU支持**：
%   ```bash
%   nvidia-smi
%   nvcc --version
%   ```
%
%
%## 九、总结
%
%VASP的编译是一项系统工程，涉及编译器选型、数学库配置、MPI并行设置以及可选的GPU加速等多个环节。本文系统梳理了VASP编译的完整知识体系：
%
%**编译器层面**，Intel oneAPI是Intel CPU平台的首选，NVIDIA HPC-SDK是GPU编译的唯一选择，AMD AOCC针对AMD EPYC优化，GCC则是开源免费的基础选项。
%
%**数学库层面**，BLAS、LAPACK、ScaLAPACK和FFTW是四大必需库。Intel MKL提供了完整的解决方案；NVIDIA HPC-SDK内置了前三者，仅需自行安装FFTW；开源方案则需要逐一配置。
%
%**MPI层面**，Intel MPI、OpenMPI和NVIDIA HPC-SDK自带的OpenMPI是三大主流选择，关键是与编译器配套使用。
%
%**GPU加速层面**，OpenACC是官方推荐的路线，自VASP 6.3.0起CUDA-C版本已被移除。编译GPU版本必须使用NVIDIA HPC-SDK（≥21.2），并注意避开问题版本22.1/22.2。
%
%**makefile.include**是编译配置的核心，VASP官方提供了丰富的模板文件，用户应根据自己的系统和需求选择合适的模板并进行定制。
%
%最后，成功的编译只是VASP使用的起点。理解编译过程的每一个环节，不仅能够帮助你顺利安装VASP，更能在后续的计算中根据需求调整编译选项，实现计算性能的最优化。正如VASP官方社区所建议的：**始终使用与VASP版本一同发布的makefile.include模板，并确保编译器、数学库和MPI库之间的兼容性**。
